\section{Dataset e preparação dos dados}
\label{sec:dataset}

O Food-101 reúne 101.000 imagens de alimentos, distribuídas igualmente entre
101 classes. A divisão oficial contém 750 imagens de treinamento e 250 de
teste por classe. As imagens de teste foram revisadas manualmente, enquanto
o conjunto de treinamento preserva ruídos, incluindo exemplos com rótulos
incorretos~\cite{bossard2014food101}. Neste trabalho, os dados são carregados
pela biblioteca Hugging Face Datasets, a partir do repositório
\texttt{ethz/food101}\footnote{\url{https://huggingface.co/datasets/ethz/food101}}.

O repositório disponibiliza dois subconjuntos, denominados \texttt{train} e
\texttt{validation}. O subconjunto original \texttt{train}, com 75.750
imagens, foi dividido em dois novos conjuntos: 80\% para treino e 20\% para
validação, com estratificação por classe e semente aleatória 42. Já o
subconjunto original \texttt{validation}, com 25.250 imagens, foi utilizado
integralmente como conjunto de teste, reservado para a avaliação final.
Portanto, a
validação utilizada durante o treinamento foi extraída exclusivamente do
\texttt{train} original. A
Tabela~\ref{tab:divisao-dados} apresenta a divisão utilizada nos quatro
experimentos. Os subconjuntos são disjuntos, e a validação é avaliada ao final
de cada época. Como a validação deriva do conjunto original de treinamento,
ela também pode conter rótulos incorretos.

\begin{table}[ht]
  \centering
  \caption{Distribuição das imagens nos subconjuntos do experimento.}
  \label{tab:divisao-dados}
  \begin{tabular}{lrr}
    \hline
    Subconjunto & Total de imagens & Imagens por classe \\
    \hline
    Treino & 60.600 & 600 \\
    Validação & 15.150 & 150 \\
    Teste & 25.250 & 250 \\
    \hline
  \end{tabular}
\end{table}

O pré-processamento é o mesmo nos três subconjuntos. As imagens são
redimensionadas para que o menor lado tenha 256 pixels, preservando a
proporção original e utilizando interpolação bicúbica. Em seguida, aplica-se
um recorte central de $224 \times 224$ pixels. Os valores dos pixels são
convertidos para ponto flutuante de 32 bits e escalados para $[0,1]$, antes
da normalização dos canais RGB com médias $(0{,}485;\,0{,}456;\,0{,}406)$ e
desvios-padrão $(0{,}229;\,0{,}224;\,0{,}225)$, adotados para a entrada do
modelo pré-treinado no ImageNet. Não são aplicadas transformações aleatórias
de aumento de dados, mantendo o pré-processamento fixo na comparação entre
as configurações do modelo.
